\documentclass[12pt]{article}

\usepackage[utf8]{inputenc}
\usepackage[T1]{fontenc}
\usepackage{geometry}
\usepackage[colorlinks=true, linkcolor=black, citecolor=black, urlcolor=blue]{hyperref}
\usepackage{amsmath}
\usepackage{graphicx}
\graphicspath{{images/}}
\usepackage{booktabs}
\usepackage{array}
\usepackage{tabularx}
\usepackage{setspace}
\usepackage{parskip}

\geometry{margin=1in}
\onehalfspacing

\title{Development of a Multimodal System for Converting PDF Documents into Structured Markdown}
\author{Noor Martini}
\date{\today}

\begin{document}

\maketitle
\tableofcontents
\newpage

% ================================================================
\section{Introduction}
% ================================================================

PDF is the dominant format for sharing and archiving documents across law,
finance, healthcare, academia and government. This is largely because a PDF
looks identical on every device and cannot be accidentally edited. But that
same property creates a serious problem for automated processing. PDF stores
text as individual characters placed at specific coordinates on a page. There
is no native concept of paragraphs, headings, or lists. Any tool that tries to extract content from
a PDF has to reconstruct that structure from scratch, using only positional
and typographic signals, which is unreliable \cite{whitington2012}.

Three main approaches exist for PDF-to-Markdown conversion. Text-based parsers
such as PyMuPDF and pdfplumber \cite{meuschke2023} read the raw text layer of
a native PDF and recover the characters, but discard layout structure in the
process. OCR tools such as Tesseract \cite{smith2007} treat the page as a flat
image and recognize characters pixel by pixel. This works for scanned documents,
but the output is plain text with no structure. A newer generation of
AI-based systems, including Marker, Microsoft MarkItDown, LlamaParse, and
multimodal models such as Mistral OCR \cite{omnidocbench2025}, uses deep learning
to recover more structure and handles tables, formulas, and multi-column layouts
more reliably. Many pipelines combine all three stages: scanned pages go through
OCR first, then a layout-aware conversion tool, then regex-based post-processing.

But all of these share one problem: they apply the same strategy to every page
of a document. A tool designed for text-heavy papers uses that same approach on
a page that contains only a scanned diagram or a dense formula. Real-world
documents mix plain text, tables, images, source code, and equations, sometimes
on the same page, and no single strategy handles all of them well. Ingemarsson
and Persson \cite{ingemarsson2024} tested this directly: rule-based tools scored
above 86\% on plain-text PDFs but dropped to 0\% on image-only pages. No single
method came out on top across all content types. And to the best of the author's
knowledge, no widely used open-source tool picks a different strategy per page
based on what that page actually contains.

This thesis builds and evaluates a system that does exactly that. Instead of
applying one fixed strategy to the whole document, the system inspects each page
individually and selects the most suitable approach based on its content type:
plain text, images, tables, hyperlinks, page numbers, source code, or
mathematical formulas. The system runs locally using LM Studio as the LLM
runtime, so there is no dependency on external cloud services and results are
fully reproducible. Converting a PDF to Markdown is more than a format change ---
it means recovering the document's structure: headings, lists, tables, links,
code blocks, and hierarchy, from a format that encodes visual appearance but not
meaning. This chapter explains the problem, defines the research questions, and
outlines the scope of the work.

% ----------------------------------------------------------------
\subsection{Problem Statement}
% ----------------------------------------------------------------

The problem this thesis addresses has two parts: the PDF format itself hides
structural information, and existing tools cannot adjust their extraction
approach based on what a given page actually contains.

\subsubsection{Difficulty of Extracting Structured Content from PDFs}

PDF was not designed to store meaning. It was designed to reproduce the
visual appearance of a printed page. The format has no native representation
of paragraphs, headings, list items, or table cells \cite{whitington2012}.
Any structure a reader sees has to be reconstructed by software from
low-level positional and typographic signals, which is fragile and often
wrong \cite{zhu2022}. A single PDF can mix plain text, embedded images,
mathematical formulas, multi-column layouts, and complex tables, sometimes
on the same page. No single extraction strategy handles all of these
correctly. Meuschke et al.\ \cite{meuschke2023} evaluated mainstream parsing
tools across multiple document domains and found that even the top performers
failed on scientific and patent documents.
The specific challenges each content type creates are covered in Chapter~2.

\subsubsection{Limited OCR Accuracy for Complex Layouts}

OCR is the standard approach for scanned PDFs, where text exists only as
pixels. But OCR accuracy drops significantly on complex layouts. Subramani
et al.\ \cite{subramani2020} describe OCR as a two-stage process: text
detection, which locates where text appears on the page, and text
transcription, which recognises characters within each detected region.
Layout analysis must happen before both stages: the tool first needs to
identify content regions (text blocks, figures, tables, equations) before
any recognition can begin. Blecher et al.\ \cite{blecher2023} and Subramani
et al.\ \cite{subramani2020} explain the mechanical limitation: traditional
OCR engines are built on a detect-first, recognise-second pattern where any
error in layout segmentation propagates directly into text recognition,
whereas modern unified visual models generate structured markup from full-page
screenshots without a separate layout step. Tools like
Tesseract~\cite{smith2007} skip this step and treat the whole page as a
single flat image, which makes it hard to assign the right semantic role to
each detected region. Zhu et al.\ \cite{zhu2022} measured the effect
directly: without layout segmentation, normalized edit distance was 0.63;
with a layout segmentation step added beforehand, it dropped to 0.38 (lower
is better). That is a large gain from one preprocessing step alone.

Multi-column layouts are a major source of errors. OCR engines often merge
text across columns or get the reading order wrong, and Zhu et
al.\ \cite{zhu2022} show that preserving correct reading order in
multi-column documents requires a dedicated segmentation stage before any
character recognition is attempted. Even current models struggle with dense
multi-column text, rotated text, tables, and mathematical notation
\cite{subramani2020, meuschke2023}. Blecher et al.\ \cite{blecher2023}
note that the PDF format causes ``a loss of semantic information,
particularly for mathematical expressions,'' and Darvishy~\cite{darvishy2018}
points out that no existing tool can interpret or describe mathematical
formulas and scientific graphs. The core issue is structural, not just
visual: what a page contains determines which extraction approach has any
chance of working.

\subsubsection{Need for Conversion to Markdown for Accessibility and Version Control}

Markdown is a lightweight markup language created by John Gruber in 2004.
It is readable as plain text and converts cleanly to HTML and other formats
\cite{cone2019}.

Converting PDFs to Markdown matters for accessibility. Darvishy
\cite{darvishy2018} explains the problem: PDF stores pages as rendering
instructions, not as structured content. Screen readers and other assistive
technologies need semantic tags and a defined reading order to process a
document correctly. Tagged PDF addresses this, but most PDFs in practice are
untagged, and Darvishy found that even when tags are present, correct reading
order and structure still require manual verification. For scientific documents
the gap is larger: no existing PDF accessibility tool can interpret
mathematical formulas or scientific graphs, which leaves that content
unreachable for screen reader users \cite{darvishy2018}. Markdown marks structure directly in the text: headings with \texttt{\#},
lists with \texttt{-}, formulas with \texttt{\$}. Assistive tools can
process it without guessing. The GitHub Flavored Markdown extension adds table syntax and
\LaTeX-style math notation, which covers most of what academic documents
require \cite{duan2025}.

Markdown's other advantages follow from being plain text. It works with
version control systems like Git, where every change is tracked at the
character level, which PDF and binary formats cannot offer. It is easy to
parse programmatically and stays readable without any specific software.
Khan et al.\ \cite{khan2024} show a concrete consequence for AI pipelines:
the quality of text extraction sets a ceiling on retrieval precision in RAG
systems. Poorly structured text produces low-quality embeddings, and retrieval
fails even when the relevant content is present. Markdown's structure prevents
that degradation.

% ----------------------------------------------------------------
\subsection{Motivation}
% ----------------------------------------------------------------

This work is motivated by a clear gap: organizations increasingly need to
process PDFs programmatically, but the tools available today either fail on
complex pages or apply the same strategy everywhere regardless of whether it
fits.

\subsubsection{Growing Importance of Document Digitization}

Document digitization is accelerating across all sectors. Organizations are
not just scanning documents anymore. They need to extract, query, and
transform document content as part of larger automated workflows. This is
driven partly by the growth of AI-based systems that consume structured text,
such as LLM-based assistants and RAG pipelines, where the quality of the
input directly affects the quality of the output \cite{khan2024}.
PDF-to-Markdown conversion is a key step in this pipeline: it bridges the
gap between static archived documents and the structured, searchable
knowledge bases that modern AI workflows depend on.

Scientific knowledge is a clear example. Blecher et al.\ \cite{blecher2023}
note that PDFs make up 2.4\% of Common Crawl and are the second most
common data format on the internet after HTML. Most of that content is
not directly machine-readable. For AI systems that need to reason over large
document collections, this is a serious bottleneck. Khan et al.\ \cite{khan2024}
tested this directly: a RAG system built on poorly extracted PDF text produced
low-quality answers even when the correct information was present in the
documents. The extraction step, not the language model, was the limiting factor.

The problem is not limited to academic papers. Ding et al.\ \cite{ding2026}
identify finance, healthcare and legal domains as sectors where visually rich
documents are the primary carrier of information. In all three, automated
workflows increasingly need to read, compare and summarise document content
at scale. PDF was not designed for this. It was designed for printing.
Closing that gap between the printed format and machine-readable output is
what motivates the work in this thesis.

\subsubsection{Advantages of Markdown as a Structured Representation}

Among common output formats, Markdown is the most practical choice for the
use cases this thesis targets. It is human-readable without a renderer,
compatible with Git for version control, easy to parse programmatically, and
well-supported by LLM and RAG tooling. A detailed comparison of PDF, Word,
and Markdown across eight relevant criteria is provided in Chapter~4.

The choice matters for downstream tasks. Khan et al.\ \cite{khan2024} show
that retrieval precision in RAG systems is directly limited by the quality of
the upstream text. Fragmented or unstructured text produces poor embeddings,
which causes retrieval to fail even when the relevant content is in the document.
Markdown preserves hierarchy through headings, lists and tables, which means
the document's structure is visible to the embedding model and to any downstream
reader. A plain paragraph dump does not provide that.

Markdown also has a practical advantage over XML and HTML for this setting.
Both of those formats require a parser to read, and neither is easy to inspect
by hand. Markdown is plain text. A developer can open the output file in any
text editor and immediately see whether the conversion was correct. For a
system that will be evaluated by comparing its output to hand-written reference
files, this readability matters: it makes the reference creation process
faster and the error analysis more direct.

\subsubsection{Comparison of Extraction Strategies}

The four strategies evaluated in this thesis differ in how they handle
different content types. Text-only extraction is fast and does not require
an LLM, but it cannot process scanned pages or pages dominated by images
or formulas. Image-only extraction sends every page through a vision LLM,
which handles all content types but applies that cost even to plain-text
pages that do not need it. The hybrid strategy improves structural fidelity
by combining both text and image signals, but still applies the same pipeline
to every page uniformly. The adaptive approach routes each page to the
strategy that fits its detected content type, which is the key difference
none of the other three provide. A full comparison across criteria such as
scanned PDF support, reading-order preservation, formula handling, and
computational cost is in Chapter~4 (Table~\ref{tab:strategy_comparison}).

Meuschke et al.\ \cite{meuschke2023} evaluated mainstream PDF parsing tools
across multiple document types and found that no single tool performed
consistently well across all of them. Rule-based text extractors performed best
on simple single-column documents but failed on scientific and patent documents.
This finding directly motivates the adaptive approach: if the best tool for one
document type is the wrong tool for another, a system that applies one strategy
uniformly will always underperform on at least some content.

Ingemarsson and Persson \cite{ingemarsson2024} reached a similar conclusion
from a different angle. They tested rule-based, deep learning-based, and
multimodal strategies on mixed-content PDFs. Rule-based tools scored 0\% on
image-only pages. Multimodal tools scored highest overall, but at a cost that
scales with every page in the document. Their results show that the right
answer depends on what is on the page, not on the document as a whole. That
is precisely what per-page adaptive routing addresses.

% ----------------------------------------------------------------
\subsection{Research Objectives and Questions}
% ----------------------------------------------------------------

The main goal of this thesis is:

\begin{quote}
\textit{To design, implement, and evaluate a multimodal, content-aware system
that converts PDF documents into structured Markdown by selecting an extraction
strategy per page based on detected content type, using LM Studio as a local
LLM runtime to ensure full reproducibility without dependence on external cloud
services, and to test whether this per-page selection outperforms any single
fixed strategy when applied to documents with diverse content types, while
reducing computational cost by skipping the LLM on plain-text pages.}
\end{quote}

In practice, plain-text pages are routed to PyMuPDF and visually complex pages
are sent to a locally hosted vision LLM via LM Studio. A single document can
use both pathways, depending on what each page contains.

This leads to five research questions.

\textbf{RQ1 --- Adaptivity.} On documents with diverse content types, does
per-page strategy selection achieve higher accuracy and structural fidelity
than any single fixed strategy applied uniformly to the full document?

\textbf{RQ2 --- Strategy comparison.} How do text-only, image-only, and hybrid
strategies compare in terms of content and structural accuracy across different
bachelor and master thesis documents?

\textbf{RQ3 --- Efficiency.} What are the computational costs (processing time,
token usage, and LLM call count) of each strategy and of the adaptive system,
and how do those costs relate to accuracy?

\textbf{RQ4 --- Output quality.} How can Markdown quality be reliably measured
beyond BLEU and ROUGE, given that structured output requires evaluating heading
hierarchy, list fidelity, and table reconstruction?

\textbf{RQ5 --- Consistency.} Do the performance differences between strategies
hold consistently across different bachelor and master thesis documents, or does
the relative ranking between strategies depend on the specific document?

These questions are tested against four hypotheses, which are evaluated against
the experimental results in Chapter~6.

\textbf{H1 --- Adaptive superiority.} On documents with diverse page content,
the adaptive system should outperform any single fixed strategy applied
uniformly, because it matches the best-suited strategy to each page type.
On documents with only one content type, adaptive should equal the
corresponding fixed strategy rather than exceed it.

\textbf{H2 --- Strategy-content fit.} Text-only will work well on plain-text
pages but fail on image-heavy or formula-rich content. Vision-based strategies
(image-only and hybrid) should handle diverse content well, which is a non-obvious result given their
black-box nature, but at significant computational cost, since every page is rendered as an image and sent through
a vision LLM regardless of whether that is actually needed.

\textbf{H3 --- Efficiency of the adaptive approach.} Despite the extra
classification step, the adaptive system should cost less on average than
image-only or hybrid. Plain-text pages skip the LLM entirely, and non-text
pages receive targeted prompts rather than a single general-purpose one,
which reduces unnecessary inference.

\textbf{H4 --- Hybrid as a strong baseline.} For documents with a balanced
mix of text and visual content, the hybrid strategy should be a competitive
baseline. The adaptive system is expected to match or beat it by avoiding
unnecessary LLM calls on plain-text pages.

% ----------------------------------------------------------------
\subsection{Scope and Delimitations}
% ----------------------------------------------------------------

This thesis works only with PDF as the input format. The system is tested
on bachelor and master thesis documents in German and English. The system
architecture is compatible with any LLM backend that supports a vision API,
but all experiments are conducted locally using LM Studio, so there is no
cloud dependency and results are fully reproducible.

The system detects the document language from the first page using the
\texttt{langdetect} library and adjusts LLM prompts to prevent the model
from translating content. If the page text is too short for reliable
detection, the system defaults to English. No language-specific training or
fine-tuning is applied, and the evaluation metrics are language-agnostic.

The following are out of scope: real-time or streaming document processing,
input formats other than PDF (e.g., DOCX, HTML, PPTX), fine-tuning or
task-specific training of LLM models, enterprise-scale deployment, cloud
infrastructure, and support for right-to-left scripts or ideographic writing
systems. The output format is Markdown only. Conversion to other structured
formats such as LaTeX, HTML, or JSON is not evaluated here, though the
architecture could support it through prompt changes.

% ----------------------------------------------------------------
\subsection{Thesis Contributions}
% ----------------------------------------------------------------

This thesis makes the following contributions.

\textbf{Per-page content classifier.} Each PDF page is classified into one
of six types (TEXT, IMAGE, FORMULA, TABLE, MIXED, or EMPTY) using structural
signals from PyMuPDF: embedded image count, vector path density, presence of
mathematical Unicode symbols, table detection, and character count. No LLM
is used at this stage, so classification is fast and deterministic.

\textbf{Font-size-based heading detection.} Heading levels are inferred from
font sizes. The system computes a character-weighted tally of font sizes on
each page to find the dominant body text size, then promotes spans that are
15\%, 40\%, and 80\% larger to Markdown heading levels H3, H2, and H1
respectively. This produces structured Markdown output for text pages
instead of flat unstructured text.

\textbf{Per-page figure extraction with Markdown linking.} Embedded figures
are extracted from vision-processed pages using PyMuPDF's
\texttt{extract\_image} API and saved as individual files (e.g.\
\texttt{figures/page\_001\_fig\_001.png}). The vision LLM is instructed to
reference these paths as Markdown image links, so figures are preserved in
the output.

\textbf{Seven-metric evaluation framework.} Output quality is measured across
seven metrics: text similarity (difflib SequenceMatcher ratio), heading
structure score, list structure score, table structure score, code block
score, paragraph structure score, and word overlap (Jaccard similarity). This
covers aspects of structural fidelity that BLEU and ROUGE cannot capture.

\textbf{Systematic experiment runner.} A configurable script tests all
combinations of strategy, model, prompt variant, temperature, and page
across the benchmark documents, records per-page scores and token usage, and
generates summary reports for strategy and model comparison.

\textbf{Reproducible open-source implementation.} The full system runs
locally using PyMuPDF, pymupdf4llm, and LM Studio. All prompt templates,
evaluation scripts, experiment configurations, and test suites are included.

% ----------------------------------------------------------------
\subsection{Thesis Structure}
% ----------------------------------------------------------------

The thesis is structured as follows. Chapter~2 provides the technical background: how PDF stores content internally, why that causes parsing problems, what Markdown is and why it suits this use case, how modern LLMs work, and how the three main extraction strategies differ. Chapter~3 reviews related work in three areas — traditional PDF tools, modern AI-based conversion systems, and LLM-based document understanding — and identifies the gaps this work addresses. Chapter~4 describes the research design, covering the document dataset, ground truth construction, evaluation metrics, and experimental setup. Chapter~5 documents the system architecture and implementation, including the per-page classifier, the four conversion strategies, the LM Studio integration, and the post-processing pipeline. Chapter~6 presents the experimental results: scores by strategy, model, temperature, and document type, with qualitative case studies of where each strategy succeeds or fails. Chapter~7 interprets those results against the five research questions and four hypotheses, examines threats to validity, and connects the findings to the broader research landscape. Chapter~8 closes with a summary of contributions, acknowledged limitations, and directions for future work.


% ================================================================
\section{Background and Fundamentals}
% ================================================================

This chapter covers the technical background needed for the rest of this thesis. It explains how PDFs are structured internally and why that creates problems for automated extraction. It also covers what Markdown is and why it works well as a target format, how modern LLMs work and what their vision capabilities mean for document conversion, and how the three main extraction strategies differ.

\subsection{PDF Format and Parsing Challenges}

PDF evolved from PostScript, a page description language built for rendering,
not for storing structured information \cite{whitington2012}. Because of this
origin, PDF encodes content as drawing instructions rather than semantic
elements. Text is stored as individual glyphs placed at $(x, y)$ coordinates
on the page, with no native concept of paragraph, heading, or list. The page
origin sits at the bottom-left corner with $Y$ increasing upward, which is the
opposite of screen coordinates. This means any tool trying to recover reading
order must first apply a coordinate transformation, and even then the result
fails on multi-column layouts \cite{whitington2012}. Any structure a reader
perceives has to be rebuilt from geometric cues alone, which is unreliable in
practice \cite{gao2011}.

The ISO 32000-1 standard introduced ``tagged PDF'' in version 1.4. Tagged PDFs
include a logical structure tree that maps rendering commands to semantic roles
such as headings and paragraphs and specifies reading order
\cite{darvishy2018}. But most PDFs in practice are untagged, so structure must
be guessed from geometry alone. And even when tags are present, correct reading
order still requires manual verification \cite{darvishy2018}. There is also an
important distinction between \textit{native PDFs}, which contain actual text
objects with bounding boxes that can be parsed directly, and \textit{scanned
PDFs}, where pages exist only as images and need OCR and visual region detection
before any text can be extracted \cite{whitington2012}. Deciding which category a
page belongs to is itself a preprocessing step, usually done by comparing the
ratio of embedded text area to image area. To make things more complicated,
some digitally produced PDFs contain visually correct but internally garbled
text encodings, which means they need an OCR fallback even though they are not
scanned \cite{whitington2012}.

Text extraction from native PDFs has its own set of problems. PDF can embed
fonts with custom encoding vectors, so character codes in the file do not
always map to Unicode. This is a common cause of garbled or incomplete
extraction, especially in older documents \cite{whitington2012}. Even when
characters are recovered correctly, simple sequential extraction still loses
paragraph boundaries, list structure, and semantic grouping
\cite{meuschke2023}. Multi-column layouts, footnotes, and mixed-language
documents make this worse, because positional order no longer corresponds to
reading order \cite{gao2011}. Formulas are a particular problem: PDF stores
them as sequences of individual characters at specific coordinates, so
superscripts, subscripts, and multi-line expressions have no semantic markup
and are often merged with body text by extractors \cite{blecher2023}.

Images embedded in a PDF carry no context linking them to the surrounding text,
so their relationship to a caption or paragraph is lost during extraction
\cite{gao2011}. Tables are similarly problematic: there is no explicit row
or column structure in PDF, and complex merged cells often require HTML output
rather than Markdown, since Markdown has no native support for \texttt{rowspan}
and \texttt{colspan} \cite{zhu2022}. Code blocks can be identified through
monospaced font detection and consistent horizontal spacing, but standard text
extractors tend to fragment them because they do not preserve
whitespace-sensitive formatting \cite{meuschke2023}.

\subsection{Markdown Format}

Markdown is a lightweight markup language created by John Gruber in 2004. It
was designed to be readable as plain text while also converting cleanly to HTML
and other formats \cite{cone2019}. The syntax uses standard ASCII characters:
\texttt{\#} symbols for headings, hyphens or asterisks for lists, backticks for
inline code, and triple backticks for fenced code blocks. Because the syntax is
so simple, a Markdown file is readable without any renderer, which is not the
case for binary or XML-based formats like PDF or DOCX.

Table~\ref{tab:markdown_syntax} lists the most commonly used Markdown syntax
elements referenced in this thesis.

\begin{table}[h!]
\centering
\caption{Common Markdown syntax elements}
\label{tab:markdown_syntax}
\renewcommand{\arraystretch}{1.4}
\begin{tabularx}{\textwidth}{llX}
\toprule
\textbf{Element} & \textbf{Syntax} & \textbf{Description} \\
\midrule
Heading 1       & \texttt{\# Heading}           & Top-level section heading \\
Heading 2       & \texttt{\#\# Heading}          & Second-level heading \\
Heading 3       & \texttt{\#\#\# Heading}        & Third-level heading \\
Bold            & \texttt{**text**}             & Bold emphasis \\
Italic          & \texttt{*text*}               & Italic emphasis \\
Unordered list  & \texttt{- item}               & Bullet list item \\
Ordered list    & \texttt{1. item}              & Numbered list item \\
Inline code     & \texttt{`code`}               & Monospace inline snippet \\
Code block      & \texttt{```lang \ldots{} ```} & Fenced block with optional language tag \\
Link            & \texttt{[text](url)}          & Hyperlink \\
Image           & \texttt{![alt](path)}         & Embedded image reference \\
Table           & \texttt{| col | col |}        & GFM pipe-delimited table \\
Math (inline)   & \texttt{\$formula\$}          & Inline \LaTeX{} expression (GFM extension) \\
\bottomrule
\end{tabularx}
\end{table}

Markdown preserves semantic hierarchy through headings, nested lists, emphasis,
and tables, which makes the output usable by both human readers and automated
tools. The GitHub Flavored Markdown (GFM) extension adds support for tables,
task lists, and \LaTeX-style math notation \cite{duan2025}, which covers most
of what academic and technical documents require. This means the output of a
PDF-to-Markdown conversion can be fed directly into documentation platforms,
static site generators, or AI pipelines without any extra transformation step.

Markdown is plain text, so it works naturally with version control systems like
Git: every change is tracked at the character level, which PDF and Word formats
cannot offer. For LLM and retrieval-augmented generation (RAG) workflows,
Markdown is a well-suited input format because it preserves document hierarchy.
Khan et al.\ \cite{khan2024} show this concretely: the quality of the upstream
text extraction sets a ceiling on downstream retrieval precision. Poorly
structured or fragmented text produces low-quality embeddings, which causes
retrieval to fail even when the relevant content is present in the document.
These are the main reasons Markdown was chosen as the target format for this
system.

\subsection{Large Language Models}

Large language models (LLMs) are neural networks trained on large text corpora
to predict and generate text. Almost all modern LLMs are built on the
transformer architecture, which uses self-attention as its core mechanism.
Self-attention lets each token in a sequence attend to every other token,
regardless of how far apart they are. This is what allows LLMs to handle long
documents and maintain coherence across a full text. Training on diverse data (books, articles, code, and web text) gives these
models broad knowledge
and the ability to perform many language tasks: translation, summarization,
question answering, and code generation, among others.

One practical constraint is the \textit{context window}: the maximum number of
tokens a model can process in a single inference call. Larger context windows
allow the model to consider longer documents without losing earlier content, but
they also require more memory and compute \cite{ingemarsson2024}. For document
conversion, the context window determines how much of a page, or how many pages at once,
can be submitted in a single prompt.

LLMs vary in capability and cost. Proprietary models like GPT-4o combine text
and visual inputs in a single end-to-end network, which lets them process
rendered page images alongside extracted text \cite{ingemarsson2024}.
Open-source alternatives like Meta's LLaMA~3 and Mistral's Mixtral reach
competitive performance at much lower inference cost, which makes them
practical for local deployment via runtimes like LM Studio. Ingemarsson and
Persson \cite{ingemarsson2024} tested several model families on PDF parsing
tasks and found that GPT-4-Turbo and GPT-4o reached the highest document
understanding scores (77--78\%), while open-source models clustered around
70\%. That gap is worth noting, but it must be weighed against the cost and
privacy benefits of running models locally.

Vision-capable LLMs are especially relevant to document conversion. By taking
rendered page images as input, they can detect visual structure that text-layer
extraction misses entirely: reading order, column boundaries, figure captions,
and mathematical notation. Prompt instructions can guide this further ---
specifying the desired output format, the document language, and the content
type of the page pushes the model toward structured, consistent output. This
combination of visual input and prompt-driven control is what the image-based
and hybrid extraction strategies in this thesis rely on.

Three parameters control how an LLM generates text and are configurable in
local runtimes such as LM Studio \cite{khan2024}. \textit{Temperature} sets
how much randomness is applied when the model picks the next token. At
temperature 0, the model always picks the most probable token, so the output
is the same every time for the same input. Higher values make the model sample
more broadly --- useful for creative tasks, but a problem for document
conversion where consistency matters. \textit{Top-K} limits candidates to the
$K$ most likely tokens at each step. \textit{Top-P} (nucleus sampling) works
differently: it selects from the smallest group of tokens whose combined
probability reaches a threshold $P$ \cite{khan2024}. At $P = 0.9$, only the
tokens that together cover 90\% of the probability mass are candidates, and
the size of that group shifts at each step based on how confident the model
is. For document conversion, temperature 0 is the right setting. It removes
randomness, ensures the same page produces the same output on every run, and
makes evaluation results repeatable.

A specific class of models has been developed for visually rich document
understanding (VRD-CU). These models treat documents as multimodal objects
where text, spatial layout, and visual features all carry meaning together.
The LayoutLM family \cite{ding2026} is a well-known example: these models are
pretrained on text content, two-dimensional positional embeddings that encode
word locations on the page, and visual features from page images. This joint
training lets the model recognize a heading not just from its words, but from
its position, font size, and surrounding whitespace --- signals that a
text-only model cannot see. The field has gone through three generations:
early feature-based models (2018--2020), BERT-style document models
(2020--2022), and LLM-instruction-tuned frameworks (2022--present)
\cite{ding2026}, which shows how quickly multimodal document understanding
has matured as a research area.

\subsection{Multimodal Processing Strategies}

PDF parsing approaches fall into three broad categories based on how they
interpret document content: rule-based text extraction, deep learning-based
recognition, and multimodal processing \cite{ingemarsson2024}. A complete
parsing system has to handle five interdependent tasks: layout analysis, text
recognition, table structure recognition, mathematical expression recognition,
and chart understanding \cite{subramani2020}. No single approach handles all
five, and the difficulty of each task varies depending on the document type.
This is one of the main reasons fixed-strategy tools struggle on diverse
real-world documents.

\textbf{Rule-based (text-only) extraction} reads the PDF's embedded text layer
directly, retrieving character sequences and their coordinates without any
visual processing. Libraries like PyMuPDF and pdfplumber work this way. They
are fast, deterministic, and reliable on native PDFs with simple layouts. But
they cannot process scanned or image-based pages at all, and they discard all
visual structure \cite{meuschke2023}. Ingemarsson and Persson
\cite{ingemarsson2024} measured this directly: rule-based tools achieve 0\%
accuracy on image-only documents, and no amount of configuration changes that.

\textbf{Deep learning-based extraction} uses trained neural network models to
recognize text, layout regions, and graphical elements from page images.
Tesseract is a well-known example; more recent tools use deeper models for
higher accuracy across a wider range of content types. These approaches can
handle scanned documents that rule-based tools cannot, and they outperform
multimodal tools on pages dominated by complex images \cite{ingemarsson2024}.
But applying them to native PDFs that already have an extractable text layer
introduces unnecessary noise and throws away existing formatting metadata
\cite{meuschke2023}. Multi-stage pipelines also have an error propagation
problem: a wrong classification in the layout analysis step affects every
stage that follows --- OCR, table detection, formula recognition --- and the
resulting errors are hard to trace back to their source \cite{duan2025}.

\textbf{Multimodal extraction} combines the text layer with visual
understanding, using models that process both the raw text content and a
rendered image of the page at the same time. These systems extract all content
types --- text, images, and tables --- process each one appropriately, and
align them according to the document's logical structure. This preserves
relationships between figures and their captions, and between equations and
the surrounding text \cite{ingemarsson2024}. Ingemarsson and Persson
\cite{ingemarsson2024} found that multimodal tools achieved the best overall
performance across diverse document types, reaching 90\% accuracy with GPT-4o.
This makes them the strongest choice when the page content is unknown or mixed.

All three strategies involve trade-offs. Rule-based methods are fast and
reliable for text but cannot handle images. Deep learning methods handle visual
content better but may over-process native PDFs. Multimodal methods offer the
broadest coverage but at higher computational cost. No single strategy is the
right choice for every page. Recent benchmarks support this: Ouyang et al.\
\cite{omnidocbench2025} benchmarked pipeline-based and end-to-end multimodal
systems across diverse document layouts in OmniDocBench and found that no single
system performs consistently well across all content types. Structural failures
persisted across all tested parsers, particularly for charts and complex visual
layouts.


% ================================================================
\section{Related Work}
% ================================================================

This chapter reviews existing work in three areas that are directly relevant to
this thesis: traditional PDF processing tools, modern document conversion
systems, and LLM-based document understanding. Each section identifies what
the existing work achieves and where the gaps are that this thesis addresses.

\subsection{Traditional PDF Processing Tools}

Traditional PDF processing tools fall into two categories: text-layer parsers
that read content directly from the embedded glyph stream of a native PDF, and
OCR-based tools that recover text by applying image recognition to rendered
page images. Both are well-established, but they were built to extract text,
not to reconstruct the semantic structure of a document. That is the core
limitation both categories share.

\subsubsection{Text-based Parsers}

Text-layer parsers like PyMuPDF, pdfplumber, and pypdfium2 read character
objects and their coordinates from a PDF's internal object stream, then
reassemble them into reading order using heuristics. Each library works at a
different level of detail. PyMuPDF gives direct access to character bounding
boxes, font metadata, and color attributes, which makes it fast and useful as
a preprocessing step. pdfplumber is built on PDFMiner and exposes bounding-box
tables for words and characters, which helps with positional grouping for
tabular data. pypdfium2 wraps the PDFium rendering engine (the same one
used in Google Chrome) and produces high-fidelity rendering, but does
little structural post-processing on its own.

Meuschke et al.\ \cite{meuschke2023} evaluated these tools across academic,
legal, and technical documents. All three libraries performed well on simple
single-column native PDFs with clean encodings, but accuracy dropped noticeably
as document complexity increased, especially on documents that mixed content
types such as embedded images, side-by-side tables, and footnotes. No single
library came out on top across all categories: PyMuPDF handled text-dense
documents fastest, while pdfplumber's coordinate data gave it an edge on
documents that required positional grouping. Ingemarsson and Persson \cite{ingemarsson2024} identified
a hard limit: on pages where content exists only as rasterised images,
text-layer parsers recover nothing, because there is no embedded text object
to read. For mathematical expressions, Gao et al.\ \cite{gao2011} show that
sequential character extraction loses the spatial relationships between glyphs
that make formulas meaningful. The output is character-correct, but the
structure is gone.

\subsubsection{OCR Solutions}

OCR tools work differently: instead of reading a text layer, they apply image
recognition to rendered page images. Smith \cite{smith2007} describes the
architecture of Tesseract, the most widely used open-source OCR engine. It
runs a two-pass pipeline: the first pass uses adaptive thresholding and
connected-component analysis to isolate character candidates, and the second
pass applies a word-level language model to resolve ambiguous boundaries.
Tesseract 4.0 added an LSTM-based recognition model that improved accuracy on
degraded and stylised text while keeping the same two-stage structure.
Subramani et al.\ \cite{subramani2020} reviewed deep learning approaches to
document understanding and placed classical engines like Tesseract in that
broader context. Their conclusion is that single-pass OCR works well on clean,
single-font documents but breaks down on complex layouts, mixed-script content,
and low-resolution scans. This is what pushed the field toward layout-aware,
multi-stage pipelines.

Zhu et al.\ \cite{zhu2022} developed DocBed, a multi-stage OCR pipeline built
for documents with complex layouts. Rather than applying recognition to the
full page at once, DocBed first runs a region segmentation model to separate
text blocks, tables, and figures, then sends each region to a dedicated
recogniser. This staged approach recovered significantly more structure than
global OCR on layout-complex documents. But it also introduced a known failure
mode: if the segmentation step makes a wrong decision, it passes incorrect
bounding boxes to every recogniser that follows, and the resulting errors are
hard to trace back to their source. This is not a problem unique to DocBed.
It applies to any modular pipeline architecture \cite{duan2025}.

\subsubsection{Shared Limitations}

Both text-layer parsers and OCR tools share a gap that Darvishy
\cite{darvishy2018} identified: no traditional tool can interpret or describe
mathematical formulas or scientific graphs in PDFs. Text-layer parsers flatten
formulas into character sequences with no spatial encoding. OCR tools recognise
individual symbols but do not rebuild the two-dimensional structure that gives
mathematical notation its meaning. Superscripts, subscripts, fraction bars,
and summation bounds are all flattened out. For the academic and scientific documents
this thesis targets, formulas and figures are not decorative. They carry
the core content, so losing their structure is a real problem. The evidence
from Meuschke et al.\ \cite{meuschke2023}, Ingemarsson and Persson
\cite{ingemarsson2024}, and Subramani et al.\ \cite{subramani2020} all point
to the same conclusion: traditional tools work well within a narrow range of
document types, but they are not sufficient for the diverse, structure-rich
documents this thesis works with. This is what motivates the move toward
learning-based and multimodal systems covered in the sections that follow.

\subsection{Modern Document Conversion Tools}

Modern conversion tools go beyond simple text extraction. Rather than reading the
raw text layer, they first analyse the visual structure of each page. This approach
is sometimes called pipeline-based document parsing: the page is decomposed into
layout blocks, each classified as text, table, figure, or heading, and the blocks
are then processed in reading order \cite{whitington2012}. Scanned pages are detected
automatically and routed to an OCR module, while digitally-born pages go through a
text extraction path. Subramani et al.\ \cite{subramani2020} survey these methods,
distinguishing between modular pipeline systems and unified Vision-Language Model
(VLM) approaches. Robust handling of complex layouts remains an open challenge
for both.

Commercial tools such as LlamaParse, Marker, and Mistral OCR follow this general
architecture but add API layers and post-processing for downstream use cases.
Ouyang et al.\ \cite{omnidocbench2025} evaluated pipeline-based and end-to-end
multimodal parsers across diverse real-world document layouts in OmniDocBench.
No single parser performed consistently well across all document types. Gaps were
largest for charts and degraded scanned pages. These results highlight how
difficult it is to build one general-purpose tool for structurally diverse
documents.

Blecher et al.\ \cite{blecher2023} introduced Nougat, a Visual Transformer
trained on paired PDF-markup data that converts full academic document pages
directly to Markdown. Nougat removes the need for a separate layout analysis
or OCR stage by generating structured output end to end. It handles headings,
equations, tables, and citations from a single model pass. Subramani et al.\
\cite{subramani2020} position this shift in a broader context: it represents
the move from pipeline-based recognition to unified visual models that reason
over the full page at once.

For long, multimodal academic documents, Xie et al.\ \cite{wukong2024} introduced
PDF-WuKong, a multimodal large language model for question-answering over long
PDFs. The system uses a sparse sampler that operates on both text and image
representations, selecting only the paragraphs and diagrams most relevant to a
given query. Evaluated on the PaperPDF dataset of 1.1 million QA pairs from
English and Chinese academic papers, PDF-WuKong outperformed proprietary products
by an average of 8.6\% on F1. This shows that efficient joint modelling of text
and images is achievable for long documents, though the system targets retrieval
and question-answering rather than full document conversion.

Even so, key limitations remain. Meuschke et al.\ \cite{meuschke2023} found
that even modern parsing tools struggle with scientific and patent documents
containing dense formulas and non-standard layouts.
Darvishy \cite{darvishy2018} notes that no existing tool provides mechanisms to
interpret and describe mathematical formulas or scientific graphs in PDFs. This gap
is one of the central motivations for the image-based extraction strategy in this
thesis, where rendered page images are passed to a multimodal LLM that can reason
about text and visual content together.

\subsection{LLM-based Document Understanding}

Before large language models, document understanding depended on chains of
specialised tools. Subramani et al.\ \cite{subramani2020} showed that even
deep-learning OCR broke down on complex layouts, mixed scripts, and
low-resolution pages. Vision-capable language models changed this: they read a
rendered page image and produce structured output token by token through
autoregressive generation, with no hand-coded layout rules in between.

Blecher et al.\ \cite{blecher2023} introduced Nougat (Neural Optical
Understanding for Academic Documents), a Visual Transformer trained on paired
PDF-markup data. It infers headings, equations, tables, and citations without a
separate layout step. Meuschke et al.\ \cite{meuschke2023} confirmed this
pattern: structured-output tools outperformed all rule-based parsers on
scientific documents. A single multimodal model can replace the multi-stage
pipeline, at least for academic documents.

OpenDataLab \cite{mineru_software} released MinerU, an open-source dual-engine
system that automatically detects whether each page is scanned or digitally born
and routes it to a Vision-Language Model or a classical OCR path. It supports
109 languages and outputs structured Markdown or JSON. Duan \cite{duan2025}
pushes the argument further:
end-to-end decoder transformers outperform pipeline-based systems on accuracy
because they never pass an incorrect intermediate result down a fixed processing
chain.

Ding et al.\ \cite{ding2026} survey Visually Rich Document Content Understanding
(VRD-CU), covering documents that combine text, layout, and visual elements.
Models pretrained across all three modalities consistently outperform text-only
approaches. Instruction-tuned LLM architectures now achieve cross-modal reasoning
that earlier models could not. Persistent gaps remain: word-level representations
miss document-level semantic structure, and multi-page documents with mixed
content are still an open challenge.

No current model covers all content types reliably. This gap motivates the
image-based strategy in this thesis, where rendered page images are passed to a
multimodal LLM and a postprocessing pipeline corrects structural failures in the
output.

\subsection{Evaluation Methods for Document Conversion}

Evaluating document conversion is harder than evaluating plain text extraction.
The output combines headings, paragraphs, tables, formulas, and figures, and each
element type can fail in a different way. Early work used string-similarity
metrics: BLEU score and edit distance (Levenshtein similarity) both compare
extracted text against a reference string, character by character
\cite{papineni2002}. Meuschke et al.\ \cite{meuschke2023} added local
alignment as a third metric, which allows gaps between matched segments and
catches word-order errors, such as merged columns or reversed reading order, that
BLEU misses.

All three metrics treat the full output as a flat sequence. A table with one
column shifted receives a low Levenshtein score even when nearly all the content
is correct. TEDS (Tree Edit Distance based Similarity) addresses this for tables:
it compares the predicted tree structure against a reference and gives partial
credit when individual cells are right even if the overall structure differs.
Research in document parsing evaluation has shown that no single metric covers
all content types, with layout analysis using mean Average Precision (mAP), text
extraction using edit metrics, and table evaluation using TEDS
\cite{zhong2020, pfitzmann2022}.

Aggregate text-similarity scores hide capability gaps. A parser may score well
under BLEU or edit distance but have a misaligned table header or a dropped
superscript that breaks any downstream process. Ouyang et al.\
\cite{omnidocbench2025} benchmarked diverse parsers across multiple content
dimensions and found that no system performed well on all types, confirming
that a single aggregate score cannot capture document parsing quality.

For conversion to Markdown specifically, N-gram metrics carry an additional
limitation: they ignore syntax. A heading marked with \texttt{\#} and the same
text written as bold have near-identical character content but different semantic
meaning in the rendered document. This is why the evaluation in Chapter 4 scores
structural, content, and formatting dimensions separately rather than collapsing
them into one number.

% ================================================================
\section{Methodology}
% ================================================================

This chapter describes the research design, dataset preparation, evaluation metrics, and experimental setup used to compare the four extraction strategies examined in this thesis.

\subsection{Research Design}

This thesis uses a comparative study to evaluate four PDF extraction strategies
against a shared test corpus and scoring framework. The comparative approach is
standard for PDF parsing research. Meuschke et al.\ \cite{meuschke2023}
applied it to benchmark PDF tools across multiple document domains. Ingemarsson
and Persson \cite{ingemarsson2024} used the same design to rank model families
by accuracy and computational cost. Using a shared methodology means the results
sit alongside that published work rather than in isolation.

The four strategies are text-only, image-only, hybrid, and adaptive. The first
three are fixed: every page goes through the same pathway regardless of content
type. The adaptive system classifies each page first, then routes it to the
appropriate extraction method. This is the only design that can skip the LLM
entirely on plain-text pages. That is what RQ3 tests.

The five research questions each target a different aspect of this comparison.
RQ1 and RQ2 compare accuracy across all four strategies: RQ1 asks whether the
adaptive system beats every fixed strategy, and RQ2 asks how the fixed
strategies rank against each other across document categories. RQ3 covers
efficiency: processing time, token usage, and LLM call count per strategy. RQ4
drives the metric design in Section~4.3. RQ5 checks whether the strategy
rankings stay consistent across different thesis documents.

Markdown is the target output format for all strategies. Unlike PDF or Word, it
stores structure in plain text, requires no proprietary software to read, and
works directly with LLM and RAG systems \cite{cone2019, khan2024}.
Table~\ref{tab:format_comparison} compares all three formats across eight
properties.

\begin{table}[h!]
\centering
\caption{Comparison of PDF, Word (.docx), and Markdown as document formats}
\label{tab:format_comparison}
\renewcommand{\arraystretch}{1.4}
\begin{tabularx}{\textwidth}{lXXX}
\toprule
\textbf{Property} & \textbf{PDF} & \textbf{Word (.docx)} & \textbf{Markdown} \\
\midrule
Human readable (raw) & No & No & \textbf{Yes} \\
Machine parseable & Difficult & Moderate & \textbf{Easy} \\
Version control (Git) & Poor & Poor & \textbf{Excellent} \\
Format portability & High & Low & \textbf{High} \\
Preserves visual fidelity & Excellent & Good & Limited \\
LLM/RAG compatibility & Poor & Moderate & \textbf{Excellent} \\
Requires proprietary software & Partial & Yes & \textbf{No} \\
Long-term archival stability & High & Moderate & \textbf{High} \\
\bottomrule
\end{tabularx}
\end{table}

\subsection{Dataset and Ground Truth Preparation}

The test corpus consists of bachelor and master thesis documents. This scope
was chosen because academic theses are structurally demanding: they combine
dense body text, mathematical formulas, multi-column layouts, citations, tables,
and figures, often within the same document. They also represent the primary
use case this system targets. Testing across multiple thesis documents rather
than a single one allows RQ5 to check whether strategy rankings are stable
across documents or vary by the specific content mix of each thesis.

Documents were drawn from real sources rather than synthetic benchmarks.
Meuschke et al.\ \cite{meuschke2023} found that tool performance varies
significantly by document structure, which is why the corpus was built to
include theses with different content mixes and layouts. The documents are
primarily in German, with a subset in English, reflecting the intended usage
context. Each selected document contains at least one structurally challenging
page: a formula-heavy section, a multi-column spread, or a table with complex
cell structure. The final corpus consists of three thesis documents totalling 90 evaluated
pages, 30 per document.

Each test document has a corresponding reference Markdown file, created by
hand. Each reference captures the full document structure: heading hierarchy,
paragraph flow, table cells, list nesting, and code blocks where present.
Pfitzmann et al.\ \cite{pfitzmann2022} argue that human-created annotation is
the only reliable baseline for evaluating document parsing tools, because
automated references embed the same systematic errors as the tools being
tested. Creating the reference by hand ensures that the evaluation measures
real output quality rather than agreement between two imperfect automated
systems.

\subsection{Evaluation Metrics}

RQ4 asks how Markdown quality can be measured beyond BLEU and ROUGE.
Section~3.4 showed that flat sequence metrics reward correct content in
the wrong structure. The answer is to score each element type separately.
Seven metrics are used, grouped into three categories: content accuracy,
structural accuracy, and formatting accuracy.

Two metrics measure how much text was preserved. Text similarity uses
difflib's SequenceMatcher, which compares candidate and reference character
by character and returns a ratio between 0 and 1. Word overlap computes
the Jaccard similarity on the word sets of both outputs. Text similarity
catches character-level noise such as OCR artefacts. Word overlap catches
word-level omissions such as missing paragraphs.

Three metrics target document structure. Heading structure score compares
the candidate and reference heading counts at each level from H1 to H6,
skipping levels that are absent in both. List structure score checks that
bullet and numbered lists appear at the correct nesting level. Paragraph
count score compares paragraph counts with a 20 per cent tolerance: a
candidate within that range of the reference scores 1.0, with a linear
penalty beyond it. A document can have high text similarity and still fail
all three of these if the extractor collapses the heading hierarchy into
plain text.

Two metrics cover formatting elements. Table detection score checks for the
presence and count of pipe tables. Code block score checks for fenced code
blocks. Both are low-frequency elements, but a single missed table is a
meaningful failure for any document that contains one. These metrics make
that failure visible rather than diluted by high text similarity on the
surrounding prose.

TEDS (Tree Edit Distance based Similarity) offers finer-grained table
evaluation by comparing cell-level tree structure \cite{zhong2020}. It was
considered but not used. The test corpus does not include tables complex
enough to require cell-level scoring, and producing cell-level annotations
by hand for a multi-document corpus carries a high annotation cost.
Table detection score covers the relevant case for this evaluation.

\subsection{Experimental Design}

The experiment applies each of the four strategies to every document in
the corpus at a fixed temperature of 0.0. Deterministic output is
required here because the evaluation compares strategies on the same
pages: any randomness introduced by higher temperatures would obscure
whether differences in score come from the strategy or from sampling
noise.

One vision-language model is used across all VLM strategies:
Qwen2.5-VL~7B, running locally through LM Studio. The text strategy
does not use a model at all. It calls PyMuPDF directly through the
pymupdf4llm library, which converts the PDF text layer to Markdown
without invoking any LLM. All four strategies therefore run under
identical hardware and model conditions, which keeps the comparison fair.

Text-only extraction at temperature 0.0 is the baseline. It is the
simplest possible extraction method: raw text from PyMuPDF with no
image input and no LLM call. A strategy must outperform it on at least
two metric categories to justify its additional computational cost.
Processing time is recorded per page, providing the raw measurements
for RQ3. Table~\ref{tab:strategy_comparison}
summarises the expected strengths and weaknesses of each approach based on
the literature \cite{meuschke2023, whitington2012}.

Results are reported as mean scores per document category across all pages
in that category. Meuschke et al.\ \cite{meuschke2023} used the same
category-level reporting and found that tools which rank well on one
document type often rank poorly on another. Category-level means make this
pattern visible. A strategy that wins overall but fails on one specific
category is a weaker result than one that performs consistently
across all four.

\begin{table}[h!]
\centering
\caption{Preliminary comparison of PDF extraction strategies}
\label{tab:strategy_comparison}
\renewcommand{\arraystretch}{1.4}
\begin{tabularx}{\textwidth}{lXXXX}
\toprule
\textbf{Criterion} & \textbf{Text-only} & \textbf{Image\nobreakdash-only} & \textbf{Hybrid} & \textbf{Adaptive} \\
\midrule
Works on scanned PDFs        & No      & Yes               & Yes          & Yes      \\
Preserves reading order      & Partial & Yes (model-dep.)  & Yes          & Yes      \\
Handles complex tables       & Poor    & Moderate          & Good         & Good     \\
Handles mathematical formulas & Poor   & Good              & Good         & Good     \\
Computational cost           & Low     & High              & \mbox{Medium-High}  & \mbox{Low-Medium} \\
Structural fidelity          & Low     & Moderate          & High         & High     \\
Dependency on LLM            & None    & Full              & Partial      & Selective \\
\bottomrule
\end{tabularx}
\end{table}


% ================================================================
\section{System Design and Implementation}
% ================================================================

This chapter describes the architecture of the developed system and the implementation of its core components: the page content classifier, the PDF processing module, the LLM integration layer, and the post-processing pipeline.

\subsection{Architecture Overview}

The system takes a PDF file as input and produces one combined Markdown document as output.
It consists of four modules: a page content classifier, a PDF processing module, an LLM integration layer, and a post-processing pipeline.
Each module is implemented as a separate Python package and communicates through well-defined data structures.
This keeps the components interchangeable and makes them testable without running the full pipeline.

Processing begins when a PDF is loaded by PyMuPDF \cite{pymupdf}.
For each page, the classifier analyses structural signals and assigns one of six content types: TEXT, IMAGE, FORMULA, TABLE, MIXED, or EMPTY.
This classification determines which extraction path the page takes.
The decision is made per page, so a single document can route different pages to different strategies.

TEXT pages are handled entirely by pymupdf4llm \cite{pymupdf4llm}, which extracts the text layer and converts it to Markdown without invoking the LLM.
All other page types are rendered as a base64-encoded PNG image and passed to a vision-capable LLM.
The LLM receives a structured prompt suited to the detected content type and returns a Markdown string.
EMPTY pages are skipped, returning a short placeholder without any processing.

Pages are converted in parallel using a thread pool. The default concurrency is four workers, which overlaps the I/O wait time of LLM calls without overloading the LM Studio server. Each worker opens its own \texttt{fitz.Document} handle, because PyMuPDF document objects are not thread-safe. The results are collected in page order and joined into a single Markdown document, with a horizontal rule between each page. The LLM integration layer communicates with LM Studio \cite{lmstudio} over a local HTTP server that exposes an OpenAI-compatible API \cite{khan2024}. All LLM calls are routed through a single \texttt{call\_llm} function, which accepts the server URL, model name, message list, temperature, and token limit as parameters. Every strategy accepts this function as an injectable parameter, so tests can replace it with a fixed-response stub without touching the rest of the code.

The post-processing pipeline receives the joined Markdown and applies a fixed sequence of cleanup passes before writing the final output file.

\subsection{Page Content Classifier}

The classifier extracts four structural signals from each page using PyMuPDF \cite{pymupdf}: the number of embedded images, the number of vector drawing paths, the total character count of the text layer, and the number of detected tables.
No machine learning is used. The classification is entirely rule-based, which keeps it fast and predictable.

Formula detection uses two independent checks.
First, the classifier counts vector paths whose bounding box is smaller than 20 pixels in either dimension.
More than 30 such short paths on a single page indicates the presence of rendered mathematical symbols, since PDF renders formula components as small vector segments.
Second, the classifier scans the text layer for mathematical Unicode characters such as $\sum$, $\int$, and $\nabla$, as well as LaTeX-style sequences that OCR engines sometimes preserve.
Either check alone is sufficient to flag the page as formula-containing.

The classification follows a fixed priority order.
A page with fewer than ten characters, no images, and fewer than five vector paths is classified as EMPTY.
A page with formula signals but fewer than 50 characters is classified as FORMULA, since formula-only pages tend to have little surrounding text.
A page with more than three embedded images is classified as IMAGE.
A page where PyMuPDF's table finder detects at least one table is classified as TABLE.
A page with any image or formula signal that did not match the earlier conditions is classified as MIXED.
A page with 50 or more characters and no visual signals is classified as TEXT.
Any remaining page falls back to MIXED with a lower confidence score of 0.5.

Each classification carries a confidence score between 0.5 and 0.95.
TEXT pages score the highest at 0.95. To reach that classification, a page must have at least 50 characters and no image, formula, or table signals. That is the hardest combination to satisfy by accident.
EMPTY and IMAGE pages score 0.9. EMPTY requires fewer than ten characters, no images, and fewer than five vector paths, which is a very narrow condition. IMAGE requires more than three embedded images, which rarely appears on a text-heavy page.
FORMULA and TABLE score 0.85 because both detectors can misfire. Vector paths above the threshold sometimes come from decorative page elements rather than formulas. PyMuPDF's table finder occasionally flags dense grid layouts that are not real tables.
MIXED pages that reached the fallback score 0.5. No single signal was strong enough, so the classification carries little certainty.
The score is stored in the \texttt{PageAnalysis} result but the current routing ignores it. A future version could use it to fall back to the hybrid strategy on low-confidence pages.

\subsection{PDF Processing Module}

The PDF processing module has two sub-pipelines: one that extracts the text layer and one that renders pages as images.
Which pipeline runs depends on the page type assigned by the classifier.
TEXT pages go through text extraction only. IMAGE, FORMULA, TABLE, and MIXED pages go through image rendering. EMPTY pages skip both.
The hybrid strategy uses both pipelines on the same page.

Text extraction is handled by pymupdf4llm \cite{pymupdf4llm}.
The system calls \texttt{pymupdf4llm.to\_markdown()} once for the whole document with \texttt{page\_chunks=True}, which returns one result per page as a list.
This matters for performance: pymupdf4llm runs a full document scan to identify header font sizes on every call, regardless of how many pages are requested.
Calling it once per page would repeat that scan N times, making extraction quadratic in the number of pages.
Calling it once and indexing into the result list keeps extraction linear.
Any inline figures found on TEXT pages are written to a configurable output directory, and their paths are corrected to be relative to the output file.

For the hybrid strategy, a second text extraction runs in parallel using a custom font-size-based extractor built on PyMuPDF \cite{pymupdf}.
This extractor builds a character-weighted tally of all font sizes on the page to identify the body text size, defined as the most frequent size by character count.
Each line's dominant font size is then compared to the body size as a ratio: lines at 1.8 times the body size or above are promoted to H1, 1.4 times to H2, and 1.15 times to H3.
This produces a lightweight Markdown representation of the text layer that the LLM receives alongside the rendered image.

Image rendering converts each page to a base64-encoded PNG using PyMuPDF's pixmap renderer.
The render resolution is 150 DPI, set via a \texttt{fitz.Matrix(150/72,~150/72)} scale factor, which is about twice the PDF native resolution of 72 DPI.
The higher resolution preserves fine detail in tables, formulas, and diagrams, which improves LLM recognition accuracy.
The encoded string is passed directly in the API request body as a data URI, so no temporary files are written to disk for vision pages.
The render DPI is a configurable constant, so it can be raised for documents with very dense content or lowered to reduce memory use.
Embedded figures can also be extracted separately using \texttt{extract\_page\_figures}, which saves each image to disk and returns relative Markdown paths that the LLM can reference in its output.
The system also preserves the printed page label for each page using PyMuPDF's \texttt{get\_label()} method, which returns the label as it appears in the PDF (for example, ``iii'' for a Roman-numeral front-matter page rather than the 0-based index). These labels appear in the \texttt{<!{-}{-} Page iii -{-}>} comments that separate pages in the output file.

\subsection{LLM Integration}

The system uses five prompt variants, each written for a specific content type.
The \textit{default} variant is a general-purpose prompt used by the text, image, and hybrid strategies when no specific variant is set.
The \textit{text} variant is used for text-classified pages and instructs the model to assign heading levels strictly by numbering depth: a top-level chapter heading becomes \texttt{\#}, a section becomes \texttt{\#\#}, a subsection \texttt{\#\#\#}, and so on.
The \textit{formula} variant instructs the model to transcribe all mathematical content as \LaTeX{} notation and to preserve subscripts, superscripts, and operators exactly without approximation.
The \textit{table} variant focuses entirely on reconstructing pipe tables from the page image, with instructions to preserve all cell values and to approximate merged cells in flat Markdown.
The \textit{diagram} variant asks the model to extract all visible text from figures and charts (titles, axis labels, legends) and to add a short description of what the diagram conveys.

The message structure sent to the API differs by strategy.
The text strategy sends a single user message containing the extracted text as a plain string, with no image content.
The image strategy sends a multipart user message containing one or more base64-encoded images as \texttt{image\_url} content blocks; no text from the PDF is included.
The hybrid strategy sends both: the extracted text is the first content block, followed by the image.
All three strategies pass the selected system prompt in a system-role message before the user turn, using the OpenAI-compatible chat completions format (see Section~\ref{sec:background:llms}).

When embedded figures have been extracted from the page in advance, a short text block is prepended to the user message listing them as Markdown image links.
This instructs the model to include those links at the appropriate positions in its output rather than describing the figures with placeholder text.
For non-English documents, a language-preservation note is appended to the system prompt by the \texttt{with\_language\_hint} function.
The note names the document language and instructs the model to keep all text in the original language.
This prevents the model from silently translating German or French content into English.

The adaptive strategy routes each page type to a specific prompt variant.
FORMULA pages use the \textit{formula} prompt, TABLE pages the \textit{table} prompt, IMAGE pages the \textit{diagram} prompt, and MIXED pages the general \textit{default} prompt.
TEXT pages are handled differently: the system calls the text strategy with \texttt{llm\_call=None}, which returns the raw pymupdf4llm output without invoking the LLM at all.
This saves the full round-trip cost on every plain-text page, which is the primary source of efficiency gain for the adaptive strategy on text-heavy documents.

\subsection{Post-processing Pipeline}

The raw output from pymupdf4llm and the LLM both contain predictable artefacts that must be removed before the Markdown is usable.
The pipeline applies a fixed sequence of passes in two stages.
The first stage runs per page via \texttt{clean\_page()}, before individual pages are joined into a document.
The second stage runs on the full joined document via \texttt{postprocess\_markdown()}.
Separating the two stages matters because some artefacts, such as duplicate running headers, only become visible once multiple pages are combined.

The first group of passes removes structural noise introduced by the PDF format.
Running headers are repeated chapter or section titles that PDF places at the top of every page; pymupdf4llm extracts these as plain-text lines, which the pipeline detects using a regex pattern and removes from the first line of each page chunk.
A second pass removes them again mid-document, where they appear as isolated paragraphs surrounded by blank lines.
Lone page-number paragraphs consisting of one to three Arabic digits or Roman numerals are also stripped.
Several further passes remove artefacts specific to LLM output.
One strips decorative horizontal rules that appear directly after a heading, which come from PDF underlines misread as Markdown separators.
Another removes prompt instruction text that some models copy verbatim into their output.
A third strips \texttt{```markdown} fences that some models place around their entire response: the fence is valid syntax for a code block but not for the top-level document.
Table-of-contents pages also receive a dedicated pass that removes dot leaders, the sequences of dots or spaces used to connect section titles to page numbers in a printed layout, which have no meaning in Markdown.

The second group targets heading quality.
pymupdf4llm often collapses all headings to the same level (\texttt{\#\#}), so the pipeline re-assigns levels based on section-number patterns: a line matching \texttt{1.1.1 Title} becomes \texttt{\#\#\#}, a line matching \texttt{1.1 Title} becomes \texttt{\#\#}, and a standalone chapter title becomes \texttt{\#}.
A separate pass merges the cases where the section number and its title were emitted as two consecutive heading lines at the same level.
A further pass handles cases where pymupdf4llm emits a bare section number with no title text, for example \texttt{\#\# 1.1} with nothing after it. The pipeline looks up the full section title in the raw PyMuPDF text for the same page and patches the heading, recovering titles that the layout extraction missed.
Headings that are single words with no recognised structural keyword (such as a person's name promoted by bold font), headings whose entire content is italic, and headings that are really numbered source-code listing lines are all demoted to plain or list text.
Title-page metadata lines promoted by bold styling are demoted to plain text by identifying everything before the first structural front-matter section (Erklärung, Abstract, Inhaltsverzeichnis).
Duplicate headings produced by mid-document running headers are removed by tracking all seen heading texts and dropping any unnumbered repetition of an already-seen numbered heading.

The remaining passes handle figures, captions, and encoding.
Bold-formatted figure caption lines produced by pymupdf4llm are converted to italic and placed directly below their image link; captions that appear before their image are swapped.
OCR text blocks extracted from embedded figures and wrapped in marker lines are converted to Markdown blockquotes.
Single non-word Unicode symbols wrapped in italic markers (such as \texttt{\_{}$\Rightarrow$\_{}}) have the markers stripped.
OCR superscript artefacts of the form \texttt{\_{}EMC\_{}[2]} are replaced with the Unicode superscript character (\texttt{EMC\textsuperscript{2}}).
LaTeX math delimiters in the \texttt{\textbackslash(...\textbackslash)} and \texttt{\textbackslash[...\textbackslash]} styles are normalised to the \texttt{\$...\$} and \texttt{\$\$...\$\$} forms used throughout the output.

\subsection{Implementation Details}

The system is implemented in Python and depends on four external libraries.
PyMuPDF \cite{pymupdf} (version 1.27) provides the core PDF interface: page analysis, text extraction, image rendering, and table detection.
pymupdf4llm \cite{pymupdf4llm} (version 1.27, matching PyMuPDF) wraps it to produce Markdown output directly from the text layer.
The \texttt{requests} library handles all HTTP communication with LM Studio.
\texttt{langdetect} identifies the primary language of each page's extracted text so the correct language hint can be appended to the system prompt.
No training or fine-tuning is required. The system runs entirely with pre-trained models loaded through LM Studio \cite{lmstudio}.

Two models are loaded in LM Studio for the experiment, as described in Section~4.4.
Qwen~3.5~9B handles text pages under the text strategy and under the adaptive strategy when it routes to the text path.
Qwen2.5-VL~7B, a vision-language model, handles all visual pages under the image-only, hybrid, and adaptive vision paths.
Both models are served through LM Studio's local HTTP server, which exposes an OpenAI-compatible \texttt{/v1/chat/completions} endpoint.
The server URL and model name are passed as parameters to every strategy function, so switching to a different model or server requires no code changes.

Reproducibility is ensured at three levels.
First, all experiments use temperature 0.0, which makes the LLM output deterministic for the same input.
Second, the library versions are pinned in the requirements file, so the extraction and rendering behaviour is identical across runs.
Third, every LLM call is routed through the injectable \texttt{call\_llm} parameter, which means the full pipeline can be exercised in tests using a fake implementation that returns fixed strings, without any dependency on a running LM Studio instance.
The evaluation framework stores all results as JSON and generates reports from that file, so results can be re-analysed without re-running the experiment.


% ================================================================
\section{Results and Evaluation}
% ================================================================

This chapter presents the results of the experiment described in Chapter~4.
All four strategies were applied to 90 pages across three thesis documents,
producing 360 evaluated page conversions with zero errors.
Results are reported per strategy across seven metrics and one efficiency
measure. Each subsection connects the findings to the relevant research
question.

% ----------------------------------------------------------------
\subsection{Overall Strategy Performance (RQ1, RQ2)}
% ----------------------------------------------------------------

Table~\ref{tab:results_overview} shows the mean scores for all four
strategies across the full corpus.

\begin{table}[h!]
\centering
\caption{Mean metric scores per strategy across 90 pages (3 documents)}
\label{tab:results_overview}
\renewcommand{\arraystretch}{1.4}
\begin{tabularx}{\textwidth}{lXXXXXXX}
\toprule
\textbf{Strategy} & \textbf{Text sim.} & \textbf{Heading} & \textbf{Lists} & \textbf{Tables} & \textbf{Code} & \textbf{Paragraphs} & \textbf{Word overlap} \\
\midrule
Text-only & \textbf{0.860} & \textbf{0.565} & 0.859 & 0.922 & \textbf{0.889} & \textbf{0.726} & 0.761 \\
Adaptive  & 0.838 & 0.542 & 0.868 & \textbf{0.956} & 0.752 & 0.636 & 0.794 \\
Image-only & 0.827 & 0.429 & \textbf{0.914} & 0.944 & 0.250 & 0.156 & \textbf{0.811} \\
Hybrid    & 0.818 & 0.348 & 0.891 & 0.944 & 0.736 & 0.597 & 0.795 \\
\bottomrule
\end{tabularx}
\end{table}

The most striking result is that text-only scores highest on five of the
seven metrics, including text similarity (0.860) and code blocks (0.889).
This is a non-obvious finding. The strategy uses no LLM and costs under
200\,ms per page. Its strength comes from the document corpus: thesis
documents consist primarily of continuous prose, and PyMuPDF's text
extraction is very accurate on such content.

The adaptive strategy scores highest on table detection (0.956), which is
0.034 points above text-only and 0.012 above image-only. This reflects its
design: pages classified as TABLE or MIXED are routed to the vision-language
model, which handles table layouts better than direct text extraction.
On the remaining metrics, adaptive is consistently second or close to
second among the four strategies.

Image-only shows a severe collapse on paragraph structure (0.156) and code
blocks (0.250). The vision-language model, when given only a rendered page
image and no extracted text, does not reproduce the exact paragraph or code
block structure from the reference. The model paraphrases and reflows the
content rather than preserving it verbatim. This makes image-only unsuitable
for documents where structural fidelity matters.

Hybrid combines image input with extracted text, which partially recovers
the paragraph and code block scores (0.597 and 0.736 respectively) compared
to image-only. However, it scores lowest on heading structure (0.348) and
text similarity (0.818). Providing both modalities does not produce the best
result on any single metric.

These results partially answer RQ1 and RQ2. The adaptive strategy does not
outperform text-only on most metrics, which contradicts hypothesis H1 for
text-heavy documents. It does outperform all strategies on table structure.
On efficiency and structural robustness across all page types, the adaptive
design shows a clear advantage, which Section~6.2 details.

% ----------------------------------------------------------------
\subsection{Efficiency (RQ3)}
% ----------------------------------------------------------------

Table~\ref{tab:efficiency} shows the mean processing time per page for
each strategy.

\begin{table}[h!]
\centering
\caption{Mean processing time per page by strategy}
\label{tab:efficiency}
\renewcommand{\arraystretch}{1.4}
\begin{tabular}{lrr}
\toprule
\textbf{Strategy} & \textbf{Avg time (ms)} & \textbf{Relative to text-only} \\
\midrule
Text-only  &    197.3 & $1\times$ \\
Adaptive   & 20{,}790.1 & $105\times$ \\
Image-only & 64{,}696.3 & $328\times$ \\
Hybrid     & 68{,}009.5 & $345\times$ \\
\bottomrule
\end{tabular}
\end{table}

Text-only is by far the fastest strategy at 197\,ms per page on average.
It makes no LLM calls, so processing time is bounded by disk read and
PyMuPDF parsing speed. The adaptive strategy averages 20{,}790\,ms per page.
This is 105 times slower than text-only but 3.1 times faster than
image-only and 3.3 times faster than hybrid. The speedup comes directly
from the page classifier: TEXT pages skip the LLM entirely, and only pages
classified as IMAGE, TABLE, FORMULA, or MIXED go through the vision model.

Image-only and hybrid are the slowest strategies at 64{,}696\,ms and
68{,}009\,ms respectively. Every page is rendered as a base64-encoded image
and sent to the vision-language model, regardless of whether the page
contains any visual content. Hybrid is slightly slower than image-only
because it additionally sends the extracted text alongside the image,
increasing the input payload on every request. This extra cost does not
translate into a consistent accuracy improvement.

These results confirm hypothesis H3: the adaptive system is significantly
cheaper than image-only or hybrid while maintaining comparable or better
accuracy. The efficiency gain is most relevant for large documents, where
the cost difference between strategies grows with page count.

% ----------------------------------------------------------------
\subsection{Per-Page Analysis and Consistency (RQ5)}
% ----------------------------------------------------------------

The per-page breakdown shows that no single strategy wins every page.
Across the 30 evaluated page positions, hybrid produces the best text
similarity on 14 pages, adaptive on 11, image-only on 5, and text-only
on 2. The variation reflects the content mix within each document.
Pages with complex table or figure layouts favour hybrid and adaptive.
Dense text pages favour text-only.

Page 6 is the lowest-scoring page across all strategies, with a best
similarity of 0.728. Inspection of the reference file shows that this
page contains a table of contents with dot-leader formatting, which none
of the strategies reproduce cleanly. This is a known limitation of both
text extraction (dot leaders appear as fragmented text) and vision
extraction (the model does not reproduce the exact list format).

The strategy rankings are consistent across the three thesis documents.
Text-only leads on text similarity in all three. Adaptive leads on table
structure in all three. This consistency addresses RQ5: the relative
performance ordering does not change across documents, even though absolute
scores vary slightly depending on content density. A document with more
formula or image pages would widen the gap between adaptive and text-only
because the classifier would route more pages to the vision model.

% ----------------------------------------------------------------
\subsection{Output Quality Metrics (RQ4)}
% ----------------------------------------------------------------

The seven-metric framework reveals failure modes that a single similarity
score would hide. Image-only scores 0.827 on text similarity, which looks
acceptable. But its paragraph structure score of 0.156 and code block score
of 0.250 expose a fundamental problem: the model reproduces most of the
words but loses the structural layout entirely. A RAG system ingesting this
output would receive the right vocabulary in the wrong format, which breaks
chunk boundaries and damages retrieval precision.

The heading structure scores for all four strategies are below 0.600. This
is a consistent weakness. The strategies differ in how they assign heading
levels: text-only derives headings from font size heuristics, while the
vision strategies rely on the model to infer structure from the visual
layout. Neither approach perfectly matches the reference heading hierarchy.
This suggests that heading level detection is the most difficult structural
task across all evaluated strategies.

Table structure is the metric where the strategies diverge most clearly.
Adaptive scores 0.956, while text-only scores 0.922 and both image-only
and hybrid score 0.944. The 0.034 gap between adaptive and text-only is
meaningful: on pages where text extraction produces malformed or missing
table markup, the vision model restores it. This is the metric where
routing to a VLM provides a measurable benefit.

Word overlap scores are close across all four strategies, ranging from
0.761 (text-only) to 0.811 (image-only). This means all strategies
recover most of the correct vocabulary. The differences between strategies
are in structure and formatting, not in which words appear in the output.


% ================================================================
\section{Discussion}
% ================================================================

\subsection{Interpretation of Results}
\begin{itemize}
    \item Why Strategies Succeed or Fail
    \item Impact of Document Characteristics
\end{itemize}

\subsection{Threats to Validity}
\begin{itemize}
    \item Internal Validity
    \item External Validity
\end{itemize}

\subsection{Practical Implications}

% TODO: mention that FORMULA pages currently use a general vision LLM, but a
% dedicated math OCR tool (e.g. pix2tex) could replace the LLM for those pages,
% reducing cost without sacrificing quality. This is a concrete future optimization.

\subsection{Comparison with Related Work}


% ================================================================
\section{Conclusion}
% ================================================================

\subsection{Summary of Findings}

\subsection{Contributions and Impact}

\subsection{Limitations of the Study}

\subsection{Future Work}
\begin{itemize}
    \item Improvements to extraction strategies
    \item Support for additional document types
    \item Integration with other tools and pipelines
    \item Exploration of newer LLM models
\end{itemize}


% ================================================================
\newpage
\begin{thebibliography}{99}

% ---------------------------------------------------------------
% Chapter 2: Background and Fundamentals
% ---------------------------------------------------------------

\bibitem{whitington2012}
Whitington, J. (2012).
\textit{PDF Explained: The ISO Standard for Document Exchange}.
O'Reilly Media.

\bibitem{darvishy2018}
Darvishy, A. (2018).
PDF accessibility: Tools and challenges.
In \textit{Proceedings of the International Conference on Computers Helping
People with Special Needs (ICCHP 2018)}, LNCS vol.\ 10896, pp.\ 113--116.
Springer.
\url{https://doi.org/10.1007/978-3-319-94277-3_20}

\bibitem{gao2011}
Gao, L., Tang, Z., Lin, X., Liu, Y., Qiu, R., \& Wang, Y. (2011).
Structure extraction from PDF-based book documents.
In \textit{Proceedings of the 11th Annual International ACM/IEEE Joint
Conference on Digital Libraries (JCDL 2011)}, pp.\ 11--20. ACM.
\url{https://dl.acm.org/doi/10.1145/1998076.1998079}

\bibitem{smith2007}
Smith, R. (2007).
An overview of the Tesseract OCR engine.
In \textit{Proceedings of the Ninth International Conference on Document
Analysis and Recognition (ICDAR 2007)}, vol.\ 2, pp.\ 629--633. IEEE.
\url{https://doi.org/10.1109/ICDAR.2007.4376991}

\bibitem{cone2019}
Cone, M. (2019).
\textit{The Markdown Guide}.
Available at: \url{https://www.markdownguide.org}

% ---------------------------------------------------------------
% Chapter 3: Related Work
% -- Sections 3.1--3.2: Traditional and Modern Document Parsing --
% ---------------------------------------------------------------

\bibitem{subramani2020}
Subramani, N., Matton, A., Greaves, M., \& Lam, A. (2022).
A survey of deep learning approaches for OCR and document understanding.
\textit{ACM Computing Surveys}, 54(6), Article 124.
\url{https://doi.org/10.1145/3495521}

\bibitem{zhu2022}
Zhu, W., Sokhandan, N., Yang, G., Martin, S., \& Sathyanarayana, S. (2022).
DocBed: A multi-stage OCR solution for documents with complex layouts.
In \textit{Proceedings of the Thirty-Sixth AAAI Conference on Artificial
Intelligence (AAAI-22)}, pp.\ 12643--12649. AAAI Press.

\bibitem{ingemarsson2024}
Ingemarsson, P., \& Persson, D. (2024).
PDF Parsing, Unveiling the Most Efficient Method.
Engineering Degree Project, Linnaeus University.

\bibitem{meuschke2023}
Meuschke, N., Jagdale, A., Spinde, T., Mitrović, J., \& Gipp, B. (2023).
A benchmark of PDF information extraction tools using a multi-task and
multi-domain evaluation framework for academic documents.
In \textit{Proceedings of the 18th International Conference on Information
(iConference 2023)}, LNCS vol.\ 13972. Springer.
\url{https://doi.org/10.1007/978-3-031-28032-0_31}

\bibitem{wukong2024}
Xie, X., Yan, H., Yin, L., Liu, Y., Ding, J., Liao, M., Liu, Y., Chen, W., \& Bai, X. (2026).
PDF-WuKong: A large multimodal model for efficient long PDF reading with
end-to-end sparse sampling.
\textit{International Journal of Computer Vision}, 134, 1--20.
\url{https://doi.org/10.1007/s11263-025-02384-5}

\bibitem{omnidocbench2025}
Ouyang, L., Qu, Y., Zhou, H., Zhu, J., Zhang, R., Lin, Q., Wang, B., Zhao, Z.,
Jiang, M., Zhao, X., Shi, J., Wu, F., Chu, P., Liu, M., Li, Z., Xu, C.,
Zhang, B., Shi, B., \& He, C. (2025).
OmniDocBench: Benchmarking diverse PDF document parsing with comprehensive
annotations.
In \textit{Proceedings of the IEEE/CVF Conference on Computer Vision and
Pattern Recognition (CVPR 2025)}.

% -- Section 3.3: LLM-based Document Understanding --

\bibitem{blecher2023}
Blecher, L., Cucurull, G., Scialom, T., \& Stojnic, R. (2023).
Nougat: Neural Optical Understanding for Academic Documents.
In \textit{Proceedings of the 2023 Conference on Empirical Methods in Natural
Language Processing (EMNLP 2023)}, pp.\ 3337--3349. ACL.
\url{https://aclanthology.org/2023.emnlp-main.199}

\bibitem{mineru_software}
OpenDataLab (2024).
\textit{MinerU: An open-source solution for precise document content extraction}
[Software].
Available at: \url{https://github.com/opendatalab/MinerU}

\bibitem{duan2025}
Duan, C. (2025).
Accelerating end-to-end PDF to Markdown conversion through assisted generation.
In \textit{Proceedings of the 30th International Conference on Applications
of Natural Language to Information Systems (NLDB 2025)},
\textit{Lecture Notes in Computer Science}, vol.\ 15836, pp.\ 34--48.
Springer Nature Switzerland.
\url{https://doi.org/10.1007/978-3-031-97141-9_3}

\bibitem{ding2026}
Ding, Y., Han, S.C., Lee, J., \& Hovy, E.H. (2026).
Deep learning based visually rich document content understanding: A survey.
\textit{Artificial Intelligence Review}, 59(4), Article 114.
Springer Nature.
\url{https://doi.org/10.1007/s10462-025-11477-3}

% -- Section 3.4: Evaluation Methods --

\bibitem{papineni2002}
Papineni, K., Roukos, S., Ward, T., \& Zhu, W. (2002).
BLEU: A method for automatic evaluation of machine translation.
In \textit{Proceedings of the 40th Annual Meeting of the Association for
Computational Linguistics (ACL 2002)}, pp.\ 311--318. ACL.

\bibitem{zhong2020}
Zhong, X., ShafieiBavani, E., \& Jimeno Yepes, A. (2020).
Image-based table recognition: Data, model, and evaluation.
In \textit{Proceedings of the 16th European Conference on Computer Vision
(ECCV 2020)}, LNCS vol.\ 12366, pp.\ 564--580. Springer.
\url{https://doi.org/10.1007/978-3-030-58589-1_34}

% ---------------------------------------------------------------
% Chapter 4: Methodology
% ---------------------------------------------------------------

\bibitem{khan2024}
Khan, A.A., Hasan, M.T., Kemell, K.K., Rasku, J., \& Abrahamsson, P. (2024).
Developing retrieval augmented generation (RAG) based LLM systems from PDFs:
An experience report.
In \textit{Proceedings of the 28th International Conference on Evaluation and
Assessment in Software Engineering (EASE 2024)}, pp.\ 136--145. ACM.
\url{https://doi.org/10.1145/3661167.3661187}

\bibitem{pfitzmann2022}
Pfitzmann, B., Auer, C., Dolfi, M., Nassar, A.S., \& Staar, P. (2022).
DocLayNet: A large human-annotated dataset for document-layout analysis.
In \textit{Proceedings of the 28th ACM SIGKDD Conference on Knowledge
Discovery and Data Mining (KDD 2022)}, pp.\ 3743--3751. ACM.
\url{https://doi.org/10.1145/3534678.3539043}

% ---------------------------------------------------------------
% Chapter 5: System Design and Implementation
% ---------------------------------------------------------------

\bibitem{pymupdf}
Artifex Software, Inc. (2024).
\textit{PyMuPDF: Python bindings for MuPDF} [Software].
Available at: \url{https://pymupdf.readthedocs.io}

\bibitem{pymupdf4llm}
Artifex Software, Inc. (2024).
\textit{pymupdf4llm: PDF to Markdown conversion for LLM ingestion} [Software].
Available at: \url{https://github.com/pymupdf/RAG}

\bibitem{lmstudio}
LM Studio (2024).
\textit{LM Studio: Run large language models locally} [Software].
Available at: \url{https://lmstudio.ai}

\end{thebibliography}


% ================================================================
\appendix

\section{Implementation Code Snippets}

\section{Complete Results Tables}

\section{Sample Documents and Outputs}

\section{Prompt Templates}

\end{document}
